\section{Question, scope, and claim discipline}
\label{sec:signed-star-scope}

The motivating phrase ``promote the signed rung to proved'' can mean three
different things, only one of which is an algorithmic theorem.

\begin{enumerate}[leftmargin=*]
\item A \emph{lower bound} says that every method in a declared class must pay
      at least a certain amount.  It does not construct a method.
\item A \emph{member upper bound} analyzes one explicit algorithm and proves
      that it pays at most that amount.
\item A \emph{matching class theorem} combines the two, thereby identifying
      the optimal scale within that class on that instance family.
\end{enumerate}

This note supplies items two and three for a promised star family.  The
algorithm is not newly invented: it is the classical optimal-SOR relaxation
specialized to the star.  The new point is the exact semantic-error analysis,
which proves that this algorithm reaches the project accuracy in accelerated
work without the logarithm required by its stronger residual certificate.

\subsection{What is proved and what is compared}

The status of every principal statement is as follows.

\begin{center}
\begin{tabular}{p{0.20\textwidth}p{0.20\textwidth}p{0.49\textwidth}}
\toprule
Statement & Status & Precise scope \\
\midrule
Optimal SOR upper bound & \textbf{Proved here} & Center-seeded
$K_{1,m}$, semantic $\varepsilon_{\rm ppr}$ error, signed coordinate
relaxations, explicit output writes. \\
Signed-class lower bound & \textbf{Proved here} & Every adaptive or randomized
dissipative one-coordinate relaxation with output equal to its maintained PPR
iterate. \\
Nonnegative-residual lower bound & \textbf{Proved here} & Every one-coordinate
method that preserves the exact residual invariant and $r\ge0$. \\
FISTA leaf-star obstruction & \textbf{Source} & Standard strongly-convex
FISTA for the $\ell_1$-regularized RPPR objective, objective-gap accuracy and
the source degree-work model. \\
ASPR path obstruction & \textbf{Imported} & Literal COLT--2023 ASPR under the
corrected objective-gap contract of the proof-owning audit note. \\
AESP--LocGD star cost & \textbf{Imported} & Literal source AESP schedule with
the batched LocGD inner method and the source PPR accuracy/work model. \\
Graph-uniform accelerated local PPR & \textbf{Open} & General undirected
graphs under the shared fully charged semantic-output contract. \\
\bottomrule
\end{tabular}
\end{center}

The comparisons are intentionally not collapsed into one theorem.  In
particular,
\(\varepsilon_{\rm ppr}\) below is the semantic degree-normalized PPR error,
whereas \(\varepsilon_{\rm obj}\) in the FISTA and ASPR sections is an RPPR
objective gap.  No equality or conversion between them is assumed.

\subsection{Promised-family versus graph-uniform results}

The upper-bound algorithm is given a star with its center as the seed.  It can
verify this promise by scanning the center and then its leaves, at additional
$O(m)$ work, which is included in the final asymptotic bound.  The proof uses
the star's exact two-dimensional invariant subspace.  It therefore does not
establish the project target
\[
  \widetilde O\!\left(
      \frac{1}{\sqrt\alpha\,\varepsilon_{\rm ppr}}
  \right)
\]
on arbitrary graphs.  A general theorem would additionally need to control
support discovery, changing local faces, repeated reads, safeguarding,
certificate queries, materialization, and output.

\subsection{Why the negative examples belong in the same note}

The star theorem shows that signed local dynamics can accelerate.  It does not
say that every insertion of momentum into a local solver is safe.  The
negative examples isolate four independent bottlenecks:

\begin{enumerate}[leftmargin=*]
\item \emph{sign restriction}: forbidding negative residuals loses the
      $1/\sqrt\alpha$ scale;
\item \emph{trajectory locality}: FISTA may touch an expensive vertex outside
      the optimum's support;
\item \emph{restart locality}: a support-safe method may repeatedly re-solve
      growing prefixes;
\item \emph{certificate mismatch}: a sufficient local residual test may be
      much stronger than the required semantic output error.
\end{enumerate}

These mechanisms are logically distinct.  A solver design that repairs only
one should not be described as resolving the others.
